\newcommand{\VarRigidTran}{{\dot{\bm{c}}_0}}
\newcommand{\VarConfig}{{\bm{v}}}

\subsection{Formulation}\label{sec:transformation}

In many cases configuration-dependent interpolations of \(\VarConfig^h\) and
\(\bm{\eta}^h\) can be interpreted as simple interpolations in a
transformation of the configuration space. Transformations are typically
used in two ways:

\begin{itemize}
\item
  When \emph{interpolation} is performed through the transformed space,
  derivatives of the transformation arise under \(\mathbb{B}^*\) 
  %in the expression for the stiffness 
  which generally cannot be removed from
  %the integrals required by the 
  the \emph{internal operators} of Box~\ref{box:fea}. 
  Such a transformation is considered \emph{element-dependent},
  because it must be accounted for in a Gauss quadrature loop when an
  element performs state-determination.
\item
  When \emph{nodal values} are taken in the transformed space, the
  transformation factors out of \(\mathbb{B}^*\) and any derivatives of
  the transformation which arise from \(\mathcal{K}_G\) can be taken
  outside of the integral inner product in the same way \(\bm{\eta}_a\)
  was in deriving \eqref{eq:def-pa}. This effect is
  \emph{element-independent} and can be performed externally through the
  use of \emph{external operators}, which will be
  introduced in this section.
\end{itemize}

In general, the effective natural bases \(\{\bm{\varphi}_{u b}\}\) can
be decomposed into an element-dependent and element-independent
contribution with the form: %
\begin{equation*} 
\bm{\varphi}_{u b}(\bm{\xi}, \{\VarConfig_c\}) = \sum_d \bar{\bm{\varphi}}_{u b}\left(\bm{\xi}; \{\VarConfig_c\}\right) \, \mathbf{a}_{b d}\left(\{\VarConfig_c\}\right)
\end{equation*}
where the contribution from operators \(\{\mathbf{a}_{b d}\}\) is
\emph{element-independent}. This independence relies on the fact that in
general
\begin{equation}
  \mathbb{B}^*\left[\bm{\varphi}^{}_{u b}\mathbf{a}\right] = \mathbb{B}^*\left[\bm{\varphi}^{}_{u b}\right]\mathbf{a}
\label{eq:B-assoc}
\end{equation} 
for any linear operator \(\mathbf{a}\)
that is constant in \(\bm{\xi}\). In this section, these ideas are
explored for two distinct types of transformations, which are identified
with the names \emph{parametric} and \emph{isometric}.

\hypertarget{sec:parametric}{%
\subsubsection{Parametric Transformations}\label{sec:parametric}}

A \emph{parametric} transformation is characterized by a map
\(g: \mathscr{C}\rightarrow\bar{\mathscr{C}}\) that takes one
configuration \((\bm{x}, \bm{\Lambda})\) to a new representation
\((\bar{\bm{x}}, \bar{\bm{\Lambda}})\) in a new configuration manifold
\(\bar{\mathscr{C}}\) where \emph{natural} variations
\(\VarConfig\in T_\chi\mathscr{C}\) and \emph{transformed} variations
\(\bar{\VarConfig} \in T_{g(\chi)}\bar{\mathscr{C}}\) are smoothly related
by the push-forward \(g_*\) and pull-back \(g^*\) at each point \(\bm{\xi}\): 
%
\begin{equation}
  \bar{\VarConfig} = g_*\VarConfig
  \qquad\text{ and }\qquad
  \VarConfig = g^* \bar{\VarConfig}.
\label{eq:uubar-parametric}
\end{equation}
In the Bubnov-Galerkin formulation, \(\bm{\eta}\) follows similarly:
\begin{equation}
\bar{\bm{\eta}} = g_*\bm{\eta}
\qquad\text{ and }\qquad
\bm{\eta} = g^* \bar{\bm{\eta}}.
\label{eq:eta-parametric}
\end{equation}
%


\paragraph{Internal Operators.}\label{internal-operators}

The simple interpolations
\(\bar{\bm{\eta}}^h\), \(\bar{\VarConfig}^h\) of the transformed variations
\(\bar{\bm{\eta}}\), \(\bar{\VarConfig}\) are represented by interpolation
bases \(\left\{\bar{\bm{\varphi}}^{}_{\eta a}\right\}\),
\(\left\{\bar{\bm{\varphi}}^{}_{u a}\right\}\) in the standard manner:
%
\begin{equation}
  \bar{\bm{\eta}}^h = \sum_a \bar{\bm{\varphi}}^{}_{\eta a}(\bm{\xi})\bar{\bm{\eta}}^{}_a
  \qquad\text{ and }\qquad
  \bar{\VarConfig}^h = \sum_b \bar{\bm{\varphi}}^{}_{u b}(\bm{\xi}) \bar{\VarConfig}^{}_b.
  \label{eq:bar-interp}
\end{equation}
%
The natural linearized variations
\(\bm{\eta}^\ell\) and \(\VarConfig^\ell\) follow immediately from
\eqref{eq:bar-interp} by \eqref{eq:eta-parametric} and
\eqref{eq:uubar-parametric}:
%
\begin{equation*} 
\bm{\eta}^\ell = \left. \frac{\mathrm{d}}{\mathrm{d}\varepsilon} g^*\bar{\bm{\eta}}^h \right|_{\varepsilon=0}
\qquad\text{ and }\qquad
\VarConfig^\ell = \left. \frac{\mathrm{d}}{\mathrm{d}\varepsilon} g^*\bar{\VarConfig}^h  \right|_{\varepsilon=0}
\end{equation*}
%
and it follows that the natural bases
\(\left\{{\bm{\varphi}}^{}_{\eta a}\right\}\),
\(\left\{{\bm{\varphi}}^{}_{u a}\right\}\) for \eqref{eq:interp} are:
%
\begin{equation}
\bm{\varphi}_{\eta a} = g^* \bar{\bm{\varphi}}_{\eta a}
\qquad\text{ and }\qquad
\bm{\varphi}_{u b} = g^* \bar{\bm{\varphi}}_{u b}.
\label{eq:para-bases}
\end{equation} 
%
Thus finite element operators for
interpolation in a transformed configuration space follow by using
\eqref{eq:para-bases} in the formulas of Box~\ref{box:fea}. 
If the Petrov-Galerkin
approach is taken, only the transformation of \eqref{eq:uubar-parametric}
needs to be accounted for in order to obtain a consistent tangent
stiffness. This furnishes the discrete strain operator:
%\(\bm{B}^\star_b\) : %
\begin{equation*} 
  \bm{B}^\star_b = \mathbb{B}^* \left[g^*\bar{\bm{\varphi}}_{u b}\right].
\end{equation*}
The pull-back \(g^*\) in this expression will generally complicate
the derivation of \(\bm{B}^\star_b\), but note that the discrete equilibrium operator
\(\bm{B}_a\) is not affected by the transformation. This is a merit of the Petrov-Galerkin
method.
An example of such a formulation is found in
\cite{jelenić1998interpolation}, where developments were made using a
similar notion of ``generalized shape functions'', which served as a
partial inspiration for the general approach developed here.
A derivation of their
element within the current framework is presented in
\cite{perez2024nonlinear}. 
Note that the discrete equilibrium operator
given by \(\mathbb{B}^*\bm{\varphi}_{\eta}\) will be identical to that
of the simple formulation by \cite{simo1985finite}. %of Section~\ref{sec:element-svq}.
Moreover, the geometric stiffness differs from 
the simple formulation only by post-multiplication with
\(g^*\).
Consequently, formulating an element with
such an interpolation generally incurs little effort
over what is
required by the simple formulation
for forming the pull-back \(g^*\), and deriving
the discrete strain operator
\(\mathbb{B}^*\left[g^*\bar{\bm{\varphi}}_{u b}\right]\).

On the other hand, an element that is formulated with the
Bubnov-Galerkin methodology must represent both \(\bm{\eta}^h\) and
\(\VarConfig^h\) identically. This is accomplished by choosing each basis
\(\bm{\varphi}_{\eta a}\) to match the form of \(\bm{\varphi}_{u b}\) .
In this case the consistent discrete equilibrium operator \(\bm{B}_a\)
is exactly the transpose of the discrete strain operator
\(\bm{B}^\star_b\). 
This generally results in a more involved expression for the 
equilibrium operator, which in turn may considerably complicate the
derivation of the geometric tangent operator \(\bm{G}_{ab}\).
A prominent example of this approach is the tour de
force initiated by \cite{cardona1988beam} and brought to fruition 
by \cite{ibrahimbegović1995computational}.


\paragraph{External Operators.}\label{sec:para-param}

It is often desirable to parameterize a finite element analysis by nodal
force values \(\{\bm{p}_a\}\) and nodal perturbations
\(\{\VarConfig_b\}\) of the configuration  
which are not necessarily sampled from the interpolated fields
\(\bar{\bm{\eta}}^h\) or \(\bar{\VarConfig}^h\). An element which
interpolates nodal unknowns
\(\bar{\VarConfig}_b\) from the transformed tangent space \(T_\chi \bar{\mathscr{C}}\) can be adapted to any
other nodal parameterization by post-multiplying the effective
interpolation \(\bm{\varphi}^{}_{u b}\) with a nodal transformation
\(\mathbf{a}_b\): 
\begin{equation}
{\bm{\varphi}}_{u b} = g^*\bar{\bm{\varphi}}_{u b} \mathbf{a}_b
\label{eq:adapt-displ}
\end{equation} 
where each nodal transformation
\(\mathbf{a}_b\) is constant over \(\bm{\xi}\), and depends on the
configuration only through the nodal parameters \(\{\bar{\VarConfig}_b\}\).
Propagating \eqref{eq:adapt-displ} through the characteristic operators
of Box~\ref{box:fea} shows that the nodal operators factor out, and the effect on
%the element forces \(\bar{\bm{p}}_a\) and 
the stiffness \({\mathbf{k}}_{a b}\) can be applied externally with the operation:
\begin{equation}
\begin{aligned}
%&\bm{p}_{a} = \bar{\bm{p}}_{a} \\
&\mathbf{k}_{ab} = \bar{\mathbf{k}}_{ab} \mathbf{a}_{b}
\end{aligned}
\label{eq:adapt-petr-para}
\end{equation}
where $\bar{\mathbf{k}}_{ab}$ is the stiffness resulting from substitution of \eqref{eq:para-bases} into the operators of Box~\ref{box:fea},
before application of the operator \(\mathbf{a}_b\).

A Bubnov-Galerkin formulation will require that nodal forces are rectified in a conjugate manner by amending
the interpolations \(\bm{\varphi}^{}_{\eta a}\) with
the nodal transformations in \eqref{eq:adapt-displ}: 
\begin{equation}
{\bm{\varphi}}^{}_{\eta a} = g^*\bar{\bm{\varphi}}^{}_{\eta a} \, \mathbf{a}_a .
\label{eq:adapt-eta}
\end{equation}
Using \eqref{eq:B-assoc} and the definition of \(\bm{p}_a\) given by \eqref{eq:def-pa},
the operator \(\mathbf{a}_a\) is factored out of the inner product as follows: %
\begin{equation*} 
\begin{aligned}
\left< g^* \bar{\bm{\varphi}}_{\eta a} \mathbf{a}_a \bm{\eta}_a,\, \mathbb{B} \bm{\sigma}\right> 
&= \left< \mathbb{B}^*[g^* \bar{\bm{\varphi}}_{\eta a}] \mathbf{a}_a \bm{\eta}_a,\, \bm{\sigma}\right> \\
&= \left< \mathbf{a}_a \bm{\eta}_a,\, \bm{B}_a \bm{\sigma}\right> \\
%\bm{\eta}_a \cdot \mathrm{a}_{ij} \mathbf{e}_j \left<\mathbf{e}_i ,\, \bm{B}_a\bm{\sigma}\right>\\
&=\bm{\eta}_a\cdot\mathbf{a}^{\mathrm{t}}_a\,\bm{p}_a .
\end{aligned}
\end{equation*}
Consequently, changes to an element with the transformation 
of \eqref{eq:adapt-eta} are equivalent to externally-applied transformations with the form:
\begin{equation} 
\begin{aligned}
\bm{p}_a    &= \mathbf{a}_a^{\mathrm{t}} \, \bar{\bm{p}}_a \\
\mathbf{k}_{ab} &= \mathbf{a}_a^{\mathrm{t}} \, \bar{\mathbf{k}}_{ab} 
  + \partial^{}_b \left[\mathbf{a}_a^{\mathrm{t}} \mid \bar{\bm{p}}_a \right]
  %+ \left.\frac{\partial}{\partial \VarConfig_b}  \mathbf{a}_a^{\mathrm{t}} \, \bar{\bm{p}}_a \right|_{\bar{\bm{p}}_a}
  %+ \partial^{}_b \left[\mathbf{a}_a^{\mathrm{t}}\right] \bar{\bm{p}}_a
\end{aligned}
\label{eq:adapt-para-bubn}
\end{equation}
where $\bar{\bm{p}}_{a}$ is the resisting force resulting from substitution of \eqref{eq:para-bases} into the operators of Box~\ref{box:fea},
before use of the operator \(\mathbf{a}_a\)
.%and \(\partial_b^{} [\mathbf{a}_a^{\mathrm{t}}] \bm{p}_a\) is the partial derivative of \(\mathbf{a}_a^{\mathrm{t}}\,\bm{p}_a\) with respect
%to parameters of node \(b\) with \(\bm{p}_a\) held constant. 
The update
for \(\mathbf{k}_{ab}\) follows equivalently either by linearizing the new
discrete residual
\(\bm{\eta}_a\cdot \mathbf{a}_a^{\mathrm{t}} \, \bm{p}_a\), or by
rederiving
the operators in Box~\ref{box:fea} using \eqref{eq:adapt-eta}.

\hypertarget{remarks04}{%
\paragraph{Remarks}\label{remarks04}}

\begin{itemize}
% \tightlist
\item
  In general, the complexity of formulations based on parametric transformations %
  is more sensitive to the interpolation of \(\bm{\eta}^h\) than \(\VarConfig^h\).
\item
  When only the interpolation of the perturbation \(\VarConfig^h\) is
  transformed and \(\bm{\eta}^h\) is still represented
  with simple  interpolations in \(T_\chi\mathscr{C}\) only the operator
  \(\bm{B}^{\star}_b\) is affected. 
  Unless special attention is paid to the interpolation of $\bm{\eta}$,
  a Petrov-Galerkin approach often furnishes nodal forces $\bm{p}_a$ that
  are not conjugate to the field sampled by the nodal parameters.
  This is a peculiarity of the Petrov-Galerkin approach and does not happen 
  in a standard variational formulation that minimizes a potential.
\end{itemize}

\hypertarget{sec:isometric}{%
\subsubsection{Isometric Transformations}\label{sec:isometric}}

Isometric transformations apply a rigid body transformation on an
element-by-element basis. In contrast to parametric transformations,
these often \emph{introduce coupling} between the variables
\(\bm{x}\) and \(\bm{\Lambda}\), and additionally couple parameters at
different nodes. The transformation map \(g\) is naturally identified by
a homogeneous translation \(\bm{c}\) and a rotation \(\bm{R}\) of the
ambient space. The transformation takes a configuration
\(\chi = (\bm{x},\bm{\Lambda})\) to the transformed representation
\(\bar{\chi} = (\bar{\bm{x}},\bar{\bm{\Lambda}})\) in the standard
fashion:
\begin{equation}
  \begin{aligned}
  \bar{\bm{x}}(\bm{\xi}) &= \bm{R}^{\mathrm{t}}\left(\bm{x}(\bm{\xi}) - \bm{c} \right) \\
  \bar{\bm{\Lambda}}(\bm{\xi})&= \bm{R}^{\mathrm{t}}\bm{\Lambda}(\bm{\xi})
  \end{aligned}
  \quad\text{ and }\quad
  \begin{aligned}
  \bm{x}(\bm{\xi}) &= \bm{R}\bar{\bm{x}}(\bm{\xi}) + \bm{c}  \\
  \bm{\Lambda}(\bm{\xi})&= \bm{R}\bar{\bm{\Lambda}}(\bm{\xi}).
  \end{aligned}
  \label{eq:tether}
\end{equation}
{
Although both \(\bm{c}\) and \(\bm{R}\)
must be constant over \(\bm{\xi}\), they may depend on the value of the
configuration variables \(\bm{x}(\bm{\xi})\) and
\(\bm{\Lambda}(\bm{\xi})\) at multiple points
\(\left\{\bm{\xi}_a\right\}\) within an element. 
% A field of three orthonormal
% directors \(\{\FrameBasis_k\}\) which will be considered stationary in
% all subsequent developments.
% Additionally %consider
% the linearly independent director fields \(\left\{\bar{\FrameBasis}_k\right\}\)
% are introduced by:
%
% \begin{equation}
% \left.\begin{aligned}
% %\FrameBasis_k &\triangleq \bm{\Lambda}\FrameBasis_k \\
% \bar{\FrameBasis}_k &\triangleq \bm{R}\FrameBasis_k \\
% %\bar{\FrameBasis}_k &\triangleq \bar{\bm{\Lambda}}\FrameBasis_k \\
% \end{aligned}\right.
% \quad\text{ implying }\qquad 
% \begin{aligned}
% %\bm{\Lambda} &= \FrameBasis_k\otimes\FrameBasis_k \\
% \bm{R}       &= \bar{\FrameBasis}_k\otimes\FrameBasis_k \\
% %\bar{\bm{\Lambda}} &= \bar{\FrameBasis}_k\otimes\FrameBasis_k \\
% \end{aligned}
% \label{eq:R-director}
% \end{equation}
% where summation is implied on repeated 
% indices.
The
transformed configuration
\(\bar{\chi}_\varepsilon = (\bar{\bm{x}}_\varepsilon, \bar{\bm{\Lambda}}_\varepsilon)\)
is related to the integral curve
\(\chi_\varepsilon = (\bm{x}_\varepsilon, \bm{\Lambda}_\varepsilon)\) by
the expressions:}

\begin{equation}
%\begin{aligned}
%\bm{x}_\varepsilon  &=  \bm{x} + \varepsilon \VarConfig_x  \\
%\bm{\Lambda}_\varepsilon &= \operatorname{Exp}\left(\varepsilon\VarConfig_{\scriptscriptstyle{\Lambda}}\right) \bm{\Lambda}
%\end{aligned}
%\quad\text{ and }\quad
\begin{aligned}
\bar{\bm{x}}_\varepsilon  &= \bm{R}^{\mathrm{t}}_\varepsilon \left(\bm{x}_\varepsilon - \bm{c}_\varepsilon \right) \\
\end{aligned}
\qquad\text{ and }\qquad
\begin{aligned}
\bar{\bm{\Lambda}}_\varepsilon &= \bm{R}^{\mathrm{t}}_\varepsilon\bm{\Lambda}_\varepsilon .
\end{aligned}
\label{eq:motion}
\end{equation}
Because the coordinate transformation %\(g\) 
depends on \(\chi\), 
both \(\bm{c}\) and \(\bm{R}\) also evolve with \(\varepsilon\).
Under this motion, the
variation \(\bar{\VarConfig}_x\) of \(\bar{\bm{x}}\) is derived from
\eqref{eq:motion} as follows: 
\begin{equation}
\begin{aligned}
\bar{\VarConfig}_x &= \frac{\mathrm{d}}{\mathrm{d}\varepsilon} \bar{\bm{x}}_\varepsilon \biggr|_{\varepsilon=0} %\\& 
= \dot{\bm{R}}^{\mathrm{t}}\left(\bm{x} + \varepsilon \VarConfig_x - \bm{c}_\varepsilon \right) + \bm{R}^{\mathrm{t}}_\varepsilon\left(\VarConfig_x - \dot{\bm{c}}_\varepsilon \right) \biggr|_{\varepsilon=0}  \\
%&= \dot{\bm{R}}^{\mathrm{t}}\left(\bm{x} - \bm{c} \right) + \bm{R}^{\mathrm{t}} \left(\VarConfig_x - \VarRigidTran \right)
\end{aligned}
\label{eq:isom-dxe}
\end{equation} 
% The initial velocity of the rotation \(\dot{\bm{R}}\) is expressed with %using
% \eqref{eq:R-director} by the sum \(\dot{\bar{\FrameBasis}}_k\otimes\FrameBasis_k\),
% after recalling that the reference basis \(\{\FrameBasis_k\}\) is
% stationary by definition. 
It is simple to show that in order for an
evolving rotation \(\bm{R}_\varepsilon\) to remain a rotation (see
definition in Section~\ref{sec:rotations}), \(\dot{\bm{R}}_\varepsilon\) must
be expressible in the form \(\bm{\Omega R}_\varepsilon\) for some
skew-symmetric \(\bm{\Omega}\) with axial vector \(\bm{\omega}\), and it follows that 
% the director
% velocities \(\dot{\bar{\FrameBasis}}_k\) are expressed in terms of an
% appropriate axial vector \(\bm{\omega}\) as follows:
% %
% \begin{equation*} 
% \dot{\bar{\FrameBasis}}_k = \bm{\omega}\times\bar{\FrameBasis}_k
% \qquad\text{ and }\qquad
% \bm{\omega}^{\times} = \dot{\bar{\FrameBasis}}_k\otimes\bar{\FrameBasis}_k
% \end{equation*}
% where the notation \(\bm{\omega}^{\times}=\bm{\Omega}\) is introduced
% by \eqref{eq:hat}. Using this in
\eqref{eq:isom-dxe} furnishes: %
\begin{equation*} 
\begin{aligned}
\bar{\VarConfig}_x &= -\bm{R}^{\mathrm{t}} \bm{\omega} \times \left(\bm{x}-\bm{c}\right)
+ \bm{R}^{\mathrm{t}}\left(\VarConfig_x - \VarRigidTran\right) \\
%&= -\bm{R}^{\mathrm{t}} \bm{\omega}^{\times}\bm{R}\bm{R}^{\mathrm{t}}  \left(\bm{x}-\bm{c}\right) + \bm{R}^{\mathrm{t}}\left(\VarConfig_x - \VarRigidTran\right) \\
%&= -\bar{\bm{\omega}}^{\times} \bm{R}^{\mathrm{t}}\left(\bm{x}-\bm{c}\right)+ \bm{R}^{\mathrm{t}}\left(\VarConfig_x - \VarRigidTran\right)
\end{aligned}
\qquad\text{ and }\qquad
\VarConfig_x = \bm{R}\bar{\VarConfig}_x + \VarRigidTran + \bm{\omega}\times (\bm{x}-\bm{c}).
\end{equation*}
with \(\VarRigidTran \triangleq \dot{\bm{c}}|_{\varepsilon=0}\). 
Variations \(\bar{\VarConfig}_{\scriptscriptstyle{\Lambda}}\) of the nonlinear configuration variable
\(\bar{\bm{\Lambda}}\) follow suit: %
\begin{equation*} 
\begin{aligned}
\bar{\VarConfig}_{\scriptscriptstyle{\Lambda}} &\triangleq \frac{\mathrm{d}}{\mathrm{d}\varepsilon}\bar{\bm{\Lambda}}_{\varepsilon}\biggr|_{\varepsilon=0} %\\ &
= \dot{\bm{R}}^{\mathrm{t}}_\varepsilon \bm{\Lambda} \biggr|_{\varepsilon=0} + \bm{R}^{\mathrm{t}} \VarConfig_{\scriptscriptstyle{\Lambda}} .%\dot{\bm{\Lambda}}_{\varepsilon} 
\end{aligned}
\end{equation*}
If variations \(\bar{\VarConfig}_{\scriptscriptstyle{\Lambda}}\) of
\(\bm{\Lambda}\) are identified by a right-invariant vector field (see
Section~\ref{sec:rotations}), then the last equation above reduces to:
\begin{equation}
\begin{aligned}
\bar{\VarConfig}_{\scriptscriptstyle{\Lambda}} = \bm{R}^{\mathrm{t}}\VarConfig_{\scriptscriptstyle{\Lambda}} - \bar{\bm{\omega}}
\end{aligned}
\quad\text{ and }\quad
\begin{aligned}
\VarConfig_{\scriptscriptstyle{\Lambda}} &= \bm{R} \bar{\VarConfig}_{\scriptscriptstyle{\Lambda}} + \bm{\omega} \\
 %&= \bm{R} \left(\bar{\VarConfig}_{\scriptscriptstyle{\Lambda}}(\bm{\xi}) + \bar{\bm{\omega}}\right)\\ 
\end{aligned}
\label{eq:rottan}
\end{equation} 
where
\(\bar{\bm{\omega}} \triangleq \bm{R}^{\mathrm{t}}\bm{\omega}\). The
additive contributions from the transformation to these variations are
collected into terms
\(\bm{\psi} \triangleq (\bm{\psi}_x, \bm{\psi}_{\scriptscriptstyle{\Lambda}})\):

\begin{equation}
\begin{aligned}
\VarConfig_x &= \bm{R}\bar{\VarConfig}_x + \bm{\psi}_x \\
\VarConfig_{\scriptscriptstyle{\Lambda}} &= \bm{R}\bar{\VarConfig}_{\scriptscriptstyle{\Lambda}} +  \bm{\psi}_{\scriptscriptstyle{\Lambda}} \\
\end{aligned}
\quad\text{ with }\quad
\begin{aligned}
\bm{\psi}_x &\triangleq  \VarRigidTran - \bm{R}\left(\bar{\bm{x}} \times \bar{\bm{\omega}}\right) \\
\bm{\psi}_{\scriptscriptstyle{\Lambda}} &\triangleq  \bm{\omega} \\
\end{aligned}
\label{eq:psi}
\end{equation} 
By a slight abuse of notation, this is written more compactly as:
%
\begin{equation}
\bar{\VarConfig} = \bm{R}^{\mathrm{t}}(\VarConfig - \bm{\psi})
\qquad\text{ and }\qquad
\VarConfig = \bm{R}\bar{\VarConfig} + \bm{\psi} .
\label{eq:global-local}
\end{equation}
where the application of the
rotation \(\bm{R}\) to \(\VarConfig\) is understood to distribute over the
couple \(\VarConfig = (\VarConfig_x , \VarConfig_{\scriptscriptstyle{\Lambda}})\) so
that
\(\bm{R}\VarConfig = (\bm{R}\VarConfig_x , \bm{R}\VarConfig_{\scriptscriptstyle{\Lambda}})\).


\hypertarget{internal-operators-1}{%
\paragraph{Internal Operators.}\label{internal-operators-1}}

The split form of the inner product in \eqref{eq:conf-pair} allows
\eqref{eq:galerkin} to be written in terms of components
\(\mathcal{G}_x\) and \(\mathcal{G}_{\scriptscriptstyle{\Lambda}}\)
defined by:

\begin{equation}
\begin{array}{rccc}
\mathcal{G}\left[\bm{\eta}\right] =& \left< \bm{\eta}_x, \, \mathbb{B}_x \bm{\sigma}\right>_x &+& \left< \bm{\eta}_{\scriptscriptstyle{\Lambda}}, \, \mathbb{B}_{\scriptscriptstyle{\Lambda}} \bm{\sigma}\right>_{\scriptscriptstyle{\Lambda}} \\ %- \left< \bm{\eta}, \, \bm{f}\right> \\
\triangleq& \mathcal{G}_x[\bm{\eta}_x] &+& \mathcal{G}_{\scriptscriptstyle{\Lambda}}[\bm{\eta}_{\scriptscriptstyle{\Lambda}}] \\ %- \left< \bm{\eta}, \, \bm{f}\right> \\
\end{array}
\label{eq:galerkin-split}
\end{equation} 
where body forces are once again dropped and 
\(\bm{\eta} = (\bm{\eta}_x, \bm{\eta}_{\scriptscriptstyle{\Lambda}})\in T^{}_\chi \mathscr{C}\).
If the Bubnov-Galerkin condition is enforced, \(\bm{\eta}\) takes the
same form as \eqref{eq:global-local}, and \(\mathcal{G}\) follows: %
\begin{equation*} 
 \begin{aligned}
 \mathcal{G}_x[\bm{\eta}] &=  \left\langle\bm{R} \bar{\bm{\eta}}_x ,\, \mathbb{B}_x \bm{\sigma} \right\rangle + \left\langle \bm{\psi}_x ,\, \mathbb{B}_x \bm{\sigma}\right\rangle \\
\end{aligned}
\qquad\text{ and }\qquad
\begin{aligned}
 \mathcal{G}_{\scriptscriptstyle{\Lambda}}[\bm{\eta}]&= \left\langle\bm{R} \bar{\bm{\eta}}_{\scriptscriptstyle{\Lambda}} ,\, \mathbb{B}_{\scriptscriptstyle{\Lambda}} \bm{\sigma} \right\rangle + \left\langle \bm{\psi}_{\scriptscriptstyle{\Lambda}} ,\, \mathbb{B}_{\scriptscriptstyle{\Lambda}} \bm{\sigma}\right\rangle .
 \end{aligned}
\end{equation*}
In writing this expression, the stress divergence
\(\mathbb{B}\bm{\sigma}\) is corresponds to evaluation on
\(\chi\). However, in practice, it is more common that the divergence is formed locally in \(\bar{\chi}\). In this case,  
invariance under isometry is invoked to write:
%
\begin{equation*} 
\begin{aligned}
\mathcal{G}_x[\bm{\eta}] &=  \left\langle\bm{R} \bar{\bm{\eta}}_x, \bm{R} \bar{\mathbb{B}}_{x} \bar{\bm{\sigma}} \right\rangle + \left\langle \bm{\psi}_x ,\, \bm{R}\bar{\mathbb{B}}_{x} \bar{\bm{\sigma}}\right\rangle \\
\end{aligned}
\qquad\text{ and }\qquad
\begin{aligned}
\mathcal{G}_{\scriptscriptstyle{\Lambda}}[\bm{\eta}]
&= \left\langle\bm{R} \bar{\bm{\eta}}_{\scriptscriptstyle{\Lambda}}, \bm{R} \bar{\mathbb{B}}_{\scriptscriptstyle{\Lambda}} \bar{\bm{\sigma}} \right\rangle 
 + \left\langle \bm{\psi}_{\scriptscriptstyle{\Lambda}} ,\, \bm{R} \bar{\mathbb{B}}_{\scriptscriptstyle{\Lambda}} \bar{\bm{\sigma}}\right\rangle .
\end{aligned}
\end{equation*}
Because \(\bm{R}\) preserves inner products by definition, it can be
removed when applied to both arguments furnishing: 
\begin{equation}
\begin{aligned}
\mathcal{G}_x\left[\bm{\eta}\right] &=  \left< \bar{\bm{\eta}}_x, \bar{\mathbb{B}}_{x} \bar{\bm{\sigma}} \right> + \left< \bar{\bm{\psi}}_x ,\, \bar{\mathbb{B}}_{x} \bar{\bm{\sigma}}\right> \\
\end{aligned}
\qquad\text{ and }\qquad
\begin{aligned}
\mathcal{G}_{\scriptscriptstyle{\Lambda}}\left[\bm{\eta}\right]&= \left< \bar{\bm{\eta}}_{\scriptscriptstyle{\Lambda}}, \, \bar{\mathbb{B}}_{\scriptscriptstyle{\Lambda}} \bar{\bm{\sigma}} \right> + \left< \bar{\bm{\psi}}_{\scriptscriptstyle{\Lambda}}, \, \bar{\mathbb{B}}_{\scriptscriptstyle{\Lambda}} \bar{\bm{\sigma}}\right>
\end{aligned}
\label{eq:galerkin-local}
\end{equation}
where
\(\bar{\bm{\psi}} \triangleq \bm{R}^{\mathrm{t}}\bm{\psi}\). The term
\(\bm{\psi}\) takes the form of a rigid body mode, and consequently,
many mechanical problems, including the model investigated in
Section~\ref{sec:cosserat}, will inherently be posed such that: %
\begin{equation*} 
\mathbb{B}^*\bm{\psi} = \mathbf{0}.
\end{equation*}
This agrees with the interpretation of \(\mathbb{B}^*\) as a virtual
strain operator as suggested by \eqref{eq:vstrain}.
When indeed \(\mathbb{B}^*\bm{\psi} = \mathbf{0}\), \eqref{eq:galerkin-local} implies %suggests
that any change resulting from an isometrically transformed interpolation can
be factored out from the internal element operators 
of Box~\ref{box:fea}.
%There is nothing from the transformation that lingers in the inner product that
%has to be accounted for inside the element state-determination.
There is no trace of the transformation in the inner product featured in the element state-determination.
This is in sharp contrast with the parametric case, where one or more of the
operators in Box~\ref{box:fea} 
%operators always
need to be re-derived because \(g_*\) either appears
%becomes trapped
under one of the differential operators (\(\mathbb{B}^*\) or
\(\partial\)), or else in %or otherwise locked under 
the inner product by dependence on \(\bm{\xi}\).
Note, however, that the transformation nonetheless  
changes the representation of the nodal variables and consequently needs to be accounted for by an
external transformation to recover the original representation.

\hypertarget{external-operators}{%
\paragraph{External Operators.}\label{external-operators}}

The external effect of the transformation stems from the transformation
of the nodal values \(\{\bm{\eta}_a\}\) and \(\{\VarConfig_b\}\) to
\(\{\bar{\bm{\eta}}_a\}\) and \(\{\bar{\VarConfig}_b\}\) by
\eqref{eq:global-local}:

\begin{equation}
\begin{aligned}
\bar{\bm{\eta}}_a &\triangleq \bar{\bm{\eta}}^h(\bm{\xi}_a) \\
&=  \bm{R}^{\mathrm{t}} \left(\bm{\eta}_a - \bm{\psi}^h(\bm{\xi}_a)\right)
\end{aligned}
\qquad\text{ and }\qquad
\begin{aligned}
\bar{\VarConfig}_b &\triangleq \bar{\VarConfig}^h(\bm{\xi}_b) \\
&= \bm{R}^{\mathrm{t}} \left(\VarConfig_b - \bm{\psi}^h(\bm{\xi}_b)\right)
\end{aligned}
\label{eq:u-augmented}
\end{equation} 
Linearizing \(\psi^h\) to form linearized variations \(\VarConfig^\ell\) and \(\bm{\eta}^\ell\) requires:
%
\begin{equation}
\VarRigidTran^{\ell} = \bm{R} \sum_b \bm{V}_b\bm{R}^{\mathrm{t}}\VarConfig_b
\qquad\text{ and }\qquad
\bar{\bm{\omega}}^{\ell} = \sum_b \bm{W}_b \bm{R}^{\mathrm{t}}\VarConfig_b
\label{eq:dpsidu}
\end{equation} 
where \(\bm{V}_b\) and \(\bm{W}_b\) denote partial derivatives with respect to nodal parameters \(\VarConfig_b\): %
\begin{equation} 
%\bm{V}_b \triangleq \frac{\partial \VarRigidTran}{\partial \VarConfig_b}
\bm{V}_b \triangleq \partial_b \VarRigidTran
\qquad\text{ and }\qquad
%\bm{W}_b \triangleq \frac{\partial \bar{\bm{\omega}}}{\partial\VarConfig_b}.
\bm{W}_b \triangleq \partial_b \bar{\bm{\omega}}.
\end{equation}
Using \eqref{eq:dpsidu} with \eqref{eq:psi} furnishes:
\begin{equation*} 
\begin{aligned}
\bm{\psi}_x^{\ell} &= \bm{R} \sum_b \bm{V}_b\bm{R}^{\mathrm{t}} \VarConfig_b - \left(\bm{x}-\bm{c}\right)\times\left(\bm{R}\sum_b \bm{W}_b \bm{R}^{\mathrm{t}}\VarConfig_b\right) \\
\end{aligned}
\qquad\text{ and }\qquad
\begin{aligned}
\bm{\psi}^{\ell}_{\scriptscriptstyle{\Lambda}} &= \bm{R}\sum_b \bm{W}_b \bm{R}^{\mathrm{t}}\VarConfig_b .
\end{aligned}
\end{equation*}
Finally, the linearization of \eqref{eq:rottan} gives for the rotation:
%
\begin{equation*} 
\begin{aligned}
\VarConfig_{\scriptscriptstyle{\Lambda}}^{\ell}(\bm{\xi})
%&= \bm{R} \bar{\VarConfig}_{\scriptscriptstyle{\Lambda}}^{\ell} + \bm{\psi}^\ell_{\scriptscriptstyle{\Lambda}} \\
&= \bm{R} \sum_b
   \bar{\bm{\varphi}}_{u b}(\bm{\xi}) \bm{R}^{\mathrm{t}} (\VarConfig_b - \bm{\omega}) 
  + \bm{R} \sum_b \bm{W}_b \bm{R}^{\mathrm{t}}\VarConfig_b .
\end{aligned}
\end{equation*}
Similarly, for the translation:
%
\begin{equation*} 
\begin{aligned}
\VarConfig_x^{\ell}(\bm{\xi})
%&= \bm{R} \bar{\VarConfig}_x^{\ell} + \VarRigidTran - \left(\bm{x} - \bm{c}\right)\times \bm{\omega} \\
&= \bm{R} \bar{\VarConfig}_x^{\ell} + \sum_b \bm{V}_b \bm{R}^{\mathrm{t}} \VarConfig_b
   - \left(\bm{R}\bar{\bm{x}}^h\right)\times \left(\bm{R} \sum_b \bm{W}_b \bm{R}^{\mathrm{t}}\VarConfig_b\right) .
%&= \bm{R} \left(\Delta\bar{\bm{x}}^h 
% + \sum_b \bar{\mathbf{V}}_b\bar{\VarConfig}_b
% - \bar{\bm{x}}^h \times \left(\sum_b \bar{\mathbf{G}}_b \bar{\VarConfig}_b\right)\right) \\
\end{aligned}
\end{equation*}
In formulations where the response is invariant under translation, the term containing
\(\bm{V}_b\) above can be dropped with no effect. In this case, using
\eqref{eq:global-local}, %and \eqref{eq:local-u} 
the linearized variations of \(\chi\) are expressed %put
in the form:
%
\begin{equation}
\VarConfig^{\ell}(\bm{\xi}) = \bm{R} \sum_b
   \bar{\bm{\varphi}}_{u b}(\bm{\xi}) \sum_c \mathbf{a}_{b c} \VarConfig_c
  + \bm{\psi}^{\ell}(\bm{\xi})
\label{eq:global-u}
\end{equation}
with
\begin{equation}
\mathbf{a}_{bc} \triangleq \left(\bm{I} \delta_{b c}- \bm{\Psi}_b \bm{W}_c\right)\bm{R}^{\mathrm{t}}
\qquad\text{ with }\quad
\bm{\Psi}_b \triangleq \left[\begin{array}{c}
-\bar{\bm{x}}(\bm{\xi}_b)^{\times} \\ \mathbf{1}
\end{array}\right]
\label{eq:projector}
\end{equation}
where \(\delta_{a b}\) is the Kronecker delta, \(\mathbf{1}\) is the identity in \(T \mathscr{C}_{\scriptscriptstyle{\Lambda}}\),
and \(\bm{I}\) the identity in \(T \mathscr{C}\).
Therefore for an isometric transformation, a consistent stiffness \(\bar{\mathbf{k}}_{a c}\) is updated as follows:
\begin{equation}
  \mathbf{k}_{ab} = \sum_{c} \bar{\mathbf{k}}_{a c} \mathbf{a}_{c b}.
  \label{eq:adapt-petr-isom}
\end{equation}
Similarly, an update for work-conjugate nodal forces \(\bar{\bm{p}}_a\) is given by:
\begin{equation}
  \begin{aligned}
  \bm{p}_a &= \sum_{b} \mathbf{a}_{b a}^{\mathrm{t}} \, \bar{\bm{p}}_{b}^{} \\
  \mathbf{k}_{ab} 
      &= \sum_{c} \mathbf{a}_{c a}^{\mathrm{t}} \, \bar{\mathbf{k}}_{cb}^{}
                   + \sum_{c} \partial_b^{} \left[\mathbf{a}_{c a}^{\mathrm{t}}\mid\bar{\bm{p}}_c^{}
                   \right]
                   %+ \mathbf{k}_g\left[\bm{p}\right]_{a b}
  \end{aligned}
  \label{eq:adapt-isom-bubn}
\end{equation}

\iffalse 
where the derivatives \(\partial_b^{} \left[\mathbf{a}_{c a}^{\mathrm{t}}\mid\bm{p}_c\right]\)
%incurred stiffness \(\mathbf{k}_g\left[\bm{p}\right]_{a b}\) is obtained by differentiation of
%the expression \(\mathbf{a}_{a c}^{\mathrm{t}}\,\bm{p}_c\) which 
take the same form as those obtained by \cite{nour-omid1991finite,lecorvec2012nonlinear}:
\begin{equation} 
\begin{aligned}
%\partial^{}_b\left[\mathbf{a}_{ab}^{\mathrm{t}}\right]\bm{p}_{b}
% \bm{G}_{a b}\left[\bm{p}\right]
\mathbf{k}_g\left[\bm{p}\right]_{a b}
&\triangleq \sum_c \partial_b^{} \left[\mathbf{a}_{c a}^{\mathrm{t}}\mid\bm{p}_c\right]
&=
- \bm{R}\bar{\bm{\Sigma}}_a \bm{W}_b \bm{R}^{\mathrm{t}}
- \bm{R} \bar{\bm{W}}^{\mathrm{t}}_a \sum_c \bar{\bm{\Sigma}}^{\bm{n}}_c \left(\bm{I}\delta_{c b} - \bm{\Psi}_c\bm{W}_b\right) \bm{R}^{\mathrm{t}} 
%
% -\bm{R}\begin{bmatrix}
% \bar{\bm{n}}^{\times}_a \bm{W}_{bx} & \bar{\bm{n}}^{\times}_a \bm{W}_{b\scriptscriptstyle{\Lambda}} \\
% \bar{\bm{m}}^{\times}_a \bm{W}_{bx} & \bar{\bm{m}}^{\times}_a \bm{W}_{b\scriptscriptstyle{\Lambda}} \\
% \end{bmatrix}\bm{R}^{\mathrm{t}}
\end{aligned}
%}
\label{eq:isom-tang}
\end{equation}
%
with %
%
\begin{equation} 
 \begin{pmatrix}
 \bar{\bm{n}}_a \\ \bar{\bm{m}}_a
 \end{pmatrix}
 \triangleq \sum_b \left(\bm{I}\delta_{ab} - \bm{W}_b^{\mathrm{t}} \bm{\Psi}_a^{\mathrm{t}}\right) \bm{p}_b,
%
 \quad 
 \bar{\bm{\Sigma}}^{\bm{n}}_a \triangleq \left[\begin{array}{cc}
 \bar{\bm{n}}_a^{\times} &
 \mathbf{0} \\
 \end{array}\right],
 \quad\text{ and }\quad
%
 \bar{\bm{\Sigma}}_a\triangleq\left[\begin{array}{c}
 \bar{\bm{n}}_a^{\times} \\
 \bar{\bm{m}}_a^{\times} \\
 \end{array}\right].
\end{equation}
Note that these derivatives are derived
without explicitly identifying a connection, and consequently, the expression 
is asymmetric away from equilibrium. However, it is shown in \cite{nour-omid1991finite}
that the expression is symmetric at equilibrium and when artificially symmetrized
quadratic convergence may still be observed in Newton-Raphson iteration in the vicinity of
a solution.
\fi

\hypertarget{sec:chain}{%
\subsubsection{Composed Transformations}}\label{sec:chain}
%
Any composition of element-independent 
transformations can be accounted for by recursive application of the appropriate
external operators to the element stiffness and resisting forces.
%
This is a consequence of the chain rule.
%$\mathbf{k}_{a b}$ % by \eqref{eq:adapt-petr-para} or \eqref{eq:adapt-petr-isom}, and optionally \eqref{eq:adapt-para-bubn} or \eqref{eq:adapt-isom-bubn} when a Bubnov-Galerkin formulation is adopted. 
The external operations of the parametric and isometric transformations
share the general structure in Box~\ref{box:external-operators} where the operators \(\mathbf{a}_a\) and \(\mathbf{a}_{ab}\) are identified by a generic linear operator \(\mathbf{a}_g\) acting on a complete set of nodal parameters.
% Application of \eqref{eq:adapt-petr-para}
% is required to maintain a consistent tangent stiffness, while \eqref{eq:adapt-para-bubn}
% is required to maintain work-conjugacy of nodal parameters.


A sequence of transformations %described
in Box~\ref{box:external-operators} can be implemented in 
a \emph{wrapper element} which presents the interface of an element to
the global solution procedure but delegates the internal state-determination to a 
\emph{wrapped} element in the existing finite element library. 
% A modular implementation of this concept extends and generalizes the framework 
% by \cite{desouza2000forcebased} with subsequent development by \cite{lecorvec2012nonlinear}.
%

\begin{ambox}[External finite element operations]{box:external-operators}
The element-independent effect of a transformation is identified by three operations:
\begin{equation*}
\begin{array}{lrl}
  %&\textbf{Symbol} &\textbf{Expression} & \textbf{Description} \\[1ex] %\hline
  \text{(i)}  & g                     & \text{External transformation} \\[2ex]
  \text{(ii)} & \mathbf{a}_g          & \text{External adaptor} \\[2ex]
  \text{(iii)}& \mathbf{k}_g\left[\bm{p}\right] 
                                      & \text{External geometric stiffness}% acting on a vector \(\bm{p}\)} \\[2ex]
\end{array}
\end{equation*}
%
Operation (i) is performed before the wrapped element is invoked,
and operations (ii) and (iii) are invoked after as an update to
the forces \(\bm{p}\) and stiffness \(\mathbf{k}\)
returned by the element:
\begin{enumerate}
\item \emph{Kinematic update}:
  \begin{equation*}
    \mathbf{k} \leftarrow {\mathbf{k}} \mathbf{a}_g
  \end{equation*}
  which takes the form of \eqref{eq:adapt-petr-para} in the parametric case, and \eqref{eq:adapt-petr-isom}
  in the isometric case.
\item \emph{Conjugate update}: 
  \begin{equation*}
    \bm{p}, \mathbf{k} \;\mapsto\;
    \begin{array}{l}
    \mathbf{a}_g^{\mathrm{t}} \, {\bm{p}}, \\
    \mathbf{a}_g^{\mathrm{t}} \, {\mathbf{k}}  + \mathbf{k}_g \left[\bm{p}\right]
    \end{array}
  \end{equation*}
  which takes the form of \eqref{eq:adapt-para-bubn} in the parametric case, and \eqref{eq:adapt-isom-bubn}
  in the isometric case.
\end{enumerate}
\end{ambox}

\iffalse
%
\begin{algorithm}[!h]
\caption{\enskip Wrapper state-determination for a sequence \(\{g_{i}\}\) of \(n_{\text{tran}}\) Box~\ref{box:external-operators} transformations.}\label{alg:push}
\begin{algorithmic}
  \Function{ElementWrapper}{$\chi$, $\VarConfig_a$, \{WrapData,ElemData\}}
  \For {Transform $g = g_{i}$ with $i=1$ to $n_{\text{tran}}$} \Comment{Forward loop over transforms}
     \State { $\left(\chi, \VarConfig_a\right)$ = \Call{$g$}{$\chi$, $\VarConfig_a$, WrapData}}
  \EndFor
  \\
  \State {$\left(\bm{p}, \mathbf{k}\right)$ = \Call{Element}{$\chi$, $\VarConfig_a$, ElemData}}
  \Comment{Invoke wrapped element}
  \\
  \For {Transform $g = g_{i}$ with $i=n_{\text{tran}}$ to $1$} \Comment{Reversed loop over transforms}
    \State {$\mathbf{k} = {\mathbf{k}}\, \mathbf{a}_g$}
    \If {Bubnov}
       \State {$\bm{p} = \mathbf{a}_g^{\mathrm{t}} \, {\bm{p}}$}
       \State {$\mathbf{k} = \mathbf{a}_g^{\mathrm{t}} \, {\mathbf{k}} + \mathbf{k}_g [\bm{p}]$}
    \EndIf
  \EndFor
  \State {\Return \(\left(\bm{p}, \mathbf{k}\right)\)}
\EndFunction
\end{algorithmic}
\end{algorithm}

Algorithm 1 outlines an implementation that follows the \emph{decorator} design
pattern \citep{mckenna1997object,scott2004software}.
For this algorithm, the wrapper state-determination consists of three phases: %a forward loop and reverse loop over the transformation sequence \(\{g\}\) which are separated by a call to the element state-determination procedure.

\begin{enumerate}

\item
  The sequence of transformations \(\{g_{i}\}\) is traversed in the \emph{forward} direction, where each \(g_i\) is applied
  to the state variables \(\chi\) and \(\VarConfig_a\).
\item
  The state-determination procedure of the wrapped element is invoked to produce the nodal forces
  \(\bm{p}\) and the stiffness \(\mathbf{k}\) in terms of the transformed state variables.
\item
  The chain is traversed in the \emph{reverse} direction
  and the nodal forces \({\bm{p}}\) and the stiffness
  \({\mathbf{k}}\) are transformed according to either the Petrov-Galerkin or the Bubnov-Galerkin rules.
\end{enumerate}
%
\fi